% =============================================================================
% Section 4: Experiments
% Paper: "The Learnable Bridge: Task Fingerprinting and Adapter Composition
%         via Structured Coupling in Low-Rank Adaptation"
% Generated: 2026-03-08
% =============================================================================

\section{Experiments}
\label{sec:experiments}

We evaluate RhombiLoRA's bridge matrices across five experimental phases,
progressing from single-adapter structural analysis through multi-task
fingerprinting to bridge-level adapter composition. All experiments use the
\texttt{rhombic} library (v0.3.0) with identical bridge extraction and
spectral analysis pipelines. Code, bridge matrices, and reproduction
scripts are available at \texttt{[repository]}.

\paragraph{Common configuration.}
Unless otherwise noted, all adapters use rank $r = 24$ with $c = 6$
channels under FCC topology, yielding $6 \times 6$ bridge matrices
(36 learnable parameters per adapter). Training uses QLoRA (4-bit
NormalFloat quantization) with AdamW ($\beta_1 = 0.9$, $\beta_2 = 0.999$),
learning rate $2 \times 10^{-4}$ with cosine schedule, batch size 4,
gradient accumulation 4, and \texttt{bf16} mixed precision. Target modules
are \texttt{q\_proj}, \texttt{k\_proj}, \texttt{v\_proj}, and
\texttt{o\_proj} across all transformer layers.

% ---------------------------------------------------------------------------
\subsection{Bridge Anatomy (1.5B)}
\label{sec:exp-anatomy-1b}

\paragraph{Setup.} We train five adapter configurations on
Qwen2.5-1.5B-Instruct (28 layers, 112 total adapters) using
Alpaca-cleaned~\cite{alpaca} (52K examples) for 2,000 steps on an NVIDIA
RTX 6000 Ada (48\,GB). The configurations are: (1)~standard LoRA $r\!=\!24$,
(2)~RhombiLoRA $r\!=\!24$ with frozen identity bridge, (3)~RhombiLoRA
$r\!=\!24$ with learnable bridge (identity initialization),
(4)~standard LoRA $r\!=\!48$ (2$\times$ parameter budget), and
(5)~RhombiLoRA $r\!=\!24$ with cubic topology ($c\!=\!3$, $3 \times 3$
bridge). Configurations 1--3 and~5 have matched trainable parameter counts
($\sim$3.27M); configuration~4 doubles this to $\sim$6.54M. Training
configurations are summarized in Table~\ref{tab:exp1-configs}.

% Table: Exp 1 configurations
\begin{table}[t]
\centering
\caption{Experiment 1 configurations (Qwen2.5-1.5B-Instruct, Alpaca-cleaned, 2K steps).}
\label{tab:exp1-configs}
\small
\begin{tabular}{@{}clcccc@{}}
\toprule
\# & Configuration & Rank & Channels & Bridge & Params \\
\midrule
1 & Standard LoRA          & 24 & 6 & Frozen $I$ & 3,268,608 \\
2 & RhombiLoRA (frozen)    & 24 & 6 & Frozen $I$ & 3,268,608 \\
3 & RhombiLoRA (learnable) & 24 & 6 & Learnable  & 3,270,624 \\
4 & Standard LoRA          & 48 & 6 & Frozen $I$ & 6,537,216 \\
5 & RhombiLoRA (cubic)     & 24 & 3 & Learnable  & 3,269,112 \\
\bottomrule
\end{tabular}
\end{table}

\paragraph{Results.}
\textbf{Identity equivalence.} Configurations 1 and~2 produce loss
trajectories within $\Delta < 0.0002$ at every checkpoint, confirming that
RhombiLoRA with a frozen identity bridge recovers standard LoRA exactly.

\textbf{Rank saturation.} Configuration~4 (rank 48, $2\times$ parameters)
converges to the same final loss as rank~24 (0.376 vs.\ 0.376 at step
2000). The training-signal bottleneck at 2K steps is not rank capacity;
the bridge provides a qualitatively different kind of capacity
(cross-channel mixing) rather than more of the same.

\textbf{Structured coupling.} Configuration~3's bridge develops monotonically
increasing algebraic connectivity (Fiedler value) over training: from
$\lambda_2 = 0.000$ at initialization to $\lambda_2 = 0.037$ at step 2000.
The primary learned behavior is diagonal self-amplification ($\sim$$+$0.009
per channel); off-diagonal cross-channel coupling is an order of magnitude
smaller ($\sim$0.001) but consistently present and growing.

\textbf{Topology matters.} The critical finding: the 6-channel FCC bridge
develops \textbf{4.6$\times$} higher Fiedler value and \textbf{5.3$\times$}
higher deviation from identity than the 3-channel cubic bridge
(Table~\ref{tab:topology-effect}), with matched parameter budgets and
identical final loss. Both topologies develop the same diagonal
self-amplification ($+$0.009), but the FCC topology discovers and exploits
more cross-channel coupling pathways. The optimizer finds that 12 edges
(FCC) offer more useful mixing structure than 6 edges (cubic)---confirming
at the neural network level the algebraic connectivity advantage
established by graph-theoretic benchmarks~\cite{rhombic-paper1}.

\begin{table}[t]
\centering
\caption{Topology effect on bridge structure (Exp 1, step 2000). Identical
parameter budget, identical loss. FCC topology develops 4.6$\times$ more
algebraic connectivity.}
\label{tab:topology-effect}
\small
\begin{tabular}{@{}lccc@{}}
\toprule
Metric & FCC (6-ch) & Cubic (3-ch) & Ratio \\
\midrule
Diagonal growth         & $+$0.0087 & $+$0.0088 & 1.0$\times$ \\
Mean $|$off-diagonal$|$ & 0.00100   & 0.00060   & 1.7$\times$ \\
Fiedler $\lambda_2$     & 0.0373    & 0.0082    & 4.6$\times$ \\
Deviation from $I$      & 0.0507    & 0.0095    & 5.3$\times$ \\
\bottomrule
\end{tabular}
\end{table}

% ---------------------------------------------------------------------------
\subsection{Bridge Anatomy at Scale (7B)}
\label{sec:exp-anatomy-7b}

\paragraph{Setup.} We train RhombiLoRA ($r\!=\!24$, FCC, learnable bridge)
on Qwen2.5-7B-Instruct (28 layers, 4 modules per layer, 112 adapters) using
Alpaca-cleaned for 10,000 steps on a local RTX 6000 Ada. Bridge matrices
are extracted at 11 checkpoints (steps 0, 1K, 2K, \ldots, 10K), yielding
1,344 bridge snapshots. We analyze the final 112 bridges for positional
structure and the full trajectory for training dynamics.

\paragraph{Per-module structure.}
Bridge coupling varies significantly by attention module type. One-way
ANOVA on final Fiedler values yields $F = 6.022$, $p = 0.0008$
(Table~\ref{tab:module-anatomy}). Q-projections develop the most coupling
(mean $\lambda_2 = 0.050$), 40\% higher than V-projections ($\lambda_2 =
0.036$). The ordering Q $>$ K $>$ V follows no trivially predicted pattern.
Notably, Q-V bridge structures are more similar than Q-K (paired $t =
2.13$, $p = 0.042$), contrary to the naive expectation that the
query--key dot product would bind Q and K more tightly. The bridge tracks
\emph{information flow} (Q queries what V supplies) rather than
\emph{similarity computation} (Q$\cdot$K).

\begin{table}[t]
\centering
\caption{Per-module-type bridge statistics (Exp 2, Qwen2.5-7B, step 10K,
$n = 28$ per module). ANOVA: $F = 6.022$, $p = 0.0008$.}
\label{tab:module-anatomy}
\small
\begin{tabular}{@{}lcccc@{}}
\toprule
Module & Fiedler ($\lambda_2$) & Deviation & Eig.\ Spread & Co/Cross \\
\midrule
q\_proj & $0.050 \pm 0.014$ & $0.273 \pm 0.070$ & $0.124 \pm 0.025$ & $1.050 \pm 0.593$ \\
k\_proj & $0.039 \pm 0.014$ & $0.155 \pm 0.034$ & $0.108 \pm 0.029$ & $1.074 \pm 0.528$ \\
o\_proj & $0.035 \pm 0.016$ & $0.225 \pm 0.099$ & $0.101 \pm 0.044$ & $1.096 \pm 0.411$ \\
v\_proj & $0.036 \pm 0.016$ & $0.142 \pm 0.064$ & $0.094 \pm 0.033$ & $1.061 \pm 0.448$ \\
\bottomrule
\end{tabular}
\end{table}

\paragraph{Depth invariance.}
In contrast to module type, layer depth has no significant effect on Fiedler
values (ANOVA $F = 0.799$, $p = 0.74$; Pearson $r = 0.03$, $p = 0.87$).
The bridge develops comparable algebraic connectivity at every depth. However,
\emph{deviation} from identity does increase with depth (ANOVA $F = 2.43$,
$p = 0.001$): later layers diverge more from the identity initialization,
suggesting that while all layers develop similar coupling strength, deeper
layers develop more complex (asymmetric or non-uniform) coupling patterns.

\paragraph{Training dynamics.}
Mean Fiedler value grows monotonically across the full trajectory: $0.000
\to 0.018$ (1K) $\to 0.040$ (10K). Early--late correlation is weak
($r = 0.099$, $p = 0.30$): adapters that develop the most coupling at step
1K do not necessarily retain that ranking at step 10K. The bridge undergoes
structural reorganization during training rather than uniform amplification
of initial structure. The coefficient of variation across final bridges is
$\text{CV} = 0.408$, with a 9$\times$ range between the most coupled adapter
(Layer~2 \texttt{q\_proj}, $\lambda_2 = 0.099$) and the least (Layer~1
\texttt{v\_proj}, $\lambda_2 = 0.011$).

\paragraph{Implication.}
The bridge is \emph{architecturally sensitive} to its position within the
attention head: module type explains significant variance in coupling
strength. This motivates the hypothesis that bridges may also be
\emph{task-sensitive}---if a 36-parameter matrix can distinguish
\texttt{q\_proj} from \texttt{v\_proj}, it may distinguish ``code'' from
``math.'' We test this directly in \S\ref{sec:exp-fingerprint}.

% ---------------------------------------------------------------------------
\subsection{Direction vs.\ Connectivity}
\label{sec:exp-direction}

\paragraph{Motivation.}
The FCC bridge topology derives from the rhombic dodecahedron, whose six
channel pairs decompose into three co-planar pairs (sharing 4 octahedral
vertices) and three cross-planar pairs (sharing 2). We hypothesized that
training data with explicit perspective-pairing structure would induce
measurably stronger co-planar than cross-planar coupling in the bridge.

\paragraph{Setup.}
We train RhombiLoRA ($r\!=\!24$, FCC) on Qwen2.5-7B-Instruct using a
synthetic geometric dataset (23K examples with 95\% co-planar perspective
pairing) for 3,000 steps on a RunPod RTX 4090. We compare against the
Alpaca baseline (Exp~2, isotropic data, 10K steps) using four spectral
diagnostics: co-planar/cross-planar coupling ratio, permutation test,
Gram matrix alignment, and perturbation retention.

\paragraph{Null result on direction.}
The geometric dataset produces \emph{no} directional signal. The co/cross
ratio is 1.002 at step 3K (effectively identity), compared to Alpaca's
transient peak of 1.091 at step 3K that decays to 1.019 by step 10K.
Neither result is statistically significant under permutation testing
(geometric $p = 0.474$; Alpaca $p = 0.332$). The full comparison is
presented in Table~\ref{tab:direction-null}.

\begin{table}[t]
\centering
\caption{Directional signal comparison (112 FCC bridges each). Geometric
data was designed with 95\% co-planar perspective pairing. Neither dataset
produces significant directional preference.}
\label{tab:direction-null}
\small
\begin{tabular}{@{}lccl@{}}
\toprule
Metric & Alpaca (10K) & Geometric (3K) & Winner \\
\midrule
Co/Cross ratio         & 1.019 & 1.002 & --- \\
Permutation $p$        & 0.332 & 0.474 & --- \\
Gram distance          & 0.921 & 0.921 & Tie \\
Eigenvalue spread      & 0.107 & 0.062 & Alpaca \\
Perturbation retention & 68.8\% & 69.7\% & Tie \\
Fiedler $\lambda_2$    & 0.040 & 0.030 & Alpaca \\
\bottomrule
\end{tabular}
\end{table}

\paragraph{Positive result on connectivity.}
While the null on direction is definitive, bridge structure is
\emph{distributed} regardless of data type: zeroing all co-planar bridge
entries retains $\sim$69\% of Fiedler value in both conditions
(68.8\% vs.\ 69.7\%). This is an architectural invariant, not a
data-dependent property. The bridge distributes its learned coupling
across both co-planar and cross-planar pairs irrespective of input
structure.

\paragraph{Root cause.}
Prompt-level co-planar bias (95\% perspective pairing) does not translate
to channel-level co-activation preference because \emph{channel assignment
is arbitrary}. The bridge has no mechanism to associate semantic perspectives
(e.g., ``analytical + empirical'') with specific channel indices (e.g.,
channels 0--1 vs.\ 2--3). Twenty-eight transformer layers between prompt
structure and bridge weights erase the prompt-level signal. Inducing
directional preference would require either explicit channel supervision,
contrastive bridge pre-training, or input-dependent bridge routing---all
of which we leave to future work.

% ---------------------------------------------------------------------------
\subsection{Task Fingerprinting}
\label{sec:exp-fingerprint}

\paragraph{Setup.}
We train three task-specific adapters on Qwen2.5-7B-Instruct:
(A)~Alpaca-cleaned (general instruction following, 52K examples, 10K steps,
local RTX 6000 Ada), (B)~CodeAlpaca-20k (code generation, 20K examples,
2K steps, local RTX 6000 Ada), and (C)~GSM8K (mathematical reasoning, 7.5K
examples, 2K steps, RunPod RTX 4090). All use identical RhombiLoRA
configuration ($r\!=\!24$, FCC, learnable bridge). Bridge matrices are
extracted at the final checkpoint, yielding 336 bridges total (112 per task,
28 layers $\times$ 4 modules). Additionally, bridges are extracted at
intermediate checkpoints (steps 0, 500, 1000, 1500) for the code and math
adapters to track fingerprint emergence.

\paragraph{Between- vs.\ within-task distance.}
Bridge distances (mean Frobenius norm across 112 positions) between tasks
are significantly larger than within-task distances. Mann-Whitney $U =
3{,}977{,}447$, $p < 10^{-6}$ (Table~\ref{tab:fingerprint-distances}).
The three task pairs are not equally distant: alpaca--code (0.189) and
alpaca--math (0.171) are well-separated, while code--math (0.066) is
substantially closer. This asymmetry is interpretable: code and math share
structured, formal reasoning patterns that distinguish both from the diverse
general-instruction domain.

\begin{table}[t]
\centering
\caption{Bridge distance statistics (Phase~1A). Between-task distances
significantly exceed within-task distances (Mann-Whitney $p < 10^{-6}$).}
\label{tab:fingerprint-distances}
\small
\begin{tabular}{@{}lrrrr@{}}
\toprule
Pair & Mean & Std & Min & Max \\
\midrule
\multicolumn{5}{@{}l}{\textit{Between-task}} \\
\quad alpaca $\leftrightarrow$ code & 0.189 & 0.087 & 0.025 & 0.443 \\
\quad alpaca $\leftrightarrow$ math & 0.171 & 0.078 & 0.039 & 0.441 \\
\quad code $\leftrightarrow$ math   & 0.066 & 0.029 & 0.024 & 0.170 \\
\midrule
\multicolumn{5}{@{}l}{\textit{Within-task}} \\
\quad alpaca & 0.177 & 0.061 & --- & --- \\
\quad code   & 0.047 & 0.015 & --- & --- \\
\quad math   & 0.085 & 0.029 & --- & --- \\
\bottomrule
\end{tabular}
\end{table}

\paragraph{Classification accuracy.}
Leave-one-out SVM classification on the 336 bridge vectors achieves 72.3\%
three-way accuracy (chance: 33.3\%). Performance is not uniform: alpaca
(93.8\%) and code (94.6\%) are well-classified, while math (50.0--66.1\%)
is frequently confused with code, consistent with their low between-task
distance.

\paragraph{Feature selection: Q-proj dominance.}
Restricting the classifier to Q-projection bridges only (28 per task, 1,008
parameters) \emph{improves} accuracy to 82.1\% (Table~\ref{tab:subset-class}).
Adding O-projection bridges (56 total, 2,016 parameters) achieves the best
overall accuracy of 83.3\%, primarily by improving math classification
from 50.0\% to 66.1\%. K-projection and V-projection bridges are
uninformative: their inclusion degrades the SVM's decision boundary.
The headline finding is that \textbf{1,008 bridge parameters (28 Q-proj
bridges) outperform the full 4,032-parameter bridge set} for task
identification. Q-projection bridges are the most task-discriminative
component of the attention head, consistent with their highest coupling
strength (\S\ref{sec:exp-anatomy-7b}).

\begin{table}[t]
\centering
\caption{Subset classification accuracy (LOO SVM, 3-way). Removing
uninformative K/V bridges improves accuracy by up to 10.4\,pp. Q-proj
bridges alone classify at 82.1\% with 1,008 parameters.}
\label{tab:subset-class}
\small
\begin{tabular}{@{}lrrr@{}}
\toprule
Subset & Positions & Parameters & LOO Accuracy \\
\midrule
Full (Q+K+V+O) & 112 & 4,032 & 72.9\% \\
Q-proj only     &  28 & 1,008 & 82.1\% \\
O-proj only     &  28 & 1,008 & 76.2\% \\
Q+O combined    &  56 & 2,016 & \textbf{83.3\%} \\
\bottomrule
\end{tabular}
\end{table}

\paragraph{Per-layer depth gradient.}
Task discrimination increases with layer depth. Spearman $\rho = 0.701$
($p < 0.0001$) between layer index and mean between-task Frobenius
distance. The top-5 most discriminative layers are all in the final quarter
of the network (layers 20--24), with a 10.3$\times$ dynamic range between
the most discriminative position (Layer~24 \texttt{o\_proj}, distance
$= 0.319$) and the least (Layer~1 \texttt{v\_proj}, distance $= 0.031$).
Late-layer O-projections surpass Q-projections at individual positions---the
output projection specializes by task in the layers where the model is
formatting its final output.

\paragraph{Magnitude, not structure.}
Tasks change bridge coupling \emph{magnitude}, not coupling
\emph{structure}. Eigenspectrum cosine similarity between task pairs is
$\geq 0.9997$ at all 112 positions. The spectral shape of the bridge
is invariant across tasks; only the scale differs. Alpaca adapters develop
the highest mean Fiedler ($0.040$), followed by math ($0.028$) and code
($0.018$). Diverse instruction-following requires more cross-channel
mixing than focused, formal domains.

\paragraph{Fingerprint emergence.}
Task fingerprints emerge entirely from training. At step~0, all bridges
are identical (identity initialization), and all between-task distances
are exactly 0.000. Divergence is monotonic: at step~1000 the binary
code--math classifier achieves 87.5\%, rising to 92.9\% at step~1500
and still climbing. The three-way Q-proj classifier reaches 84.5\% at
the final checkpoint. This trajectory rules out any initialization
confound and demonstrates that bridge structure is a faithful record
of training-induced specialization.

% ---------------------------------------------------------------------------
\subsection{Bridge-Level Adapter Merging}
\label{sec:exp-merging}

\paragraph{Setup.}
We interpolate final bridge matrices between all three task pairs at 11
uniformly spaced mixing coefficients ($\alpha \in \{0.0, 0.1, \ldots,
1.0\}$), where $B_\alpha = (1 - \alpha) B_A + \alpha B_B$. For each
$\alpha$, we compute Fiedler value, deviation from identity, Frobenius
distance to each source, and eigenspectrum cosine similarity to each
source across all 112 adapter positions. We compare against a random
baseline: merging each task bridge with 20 independent random $6 \times
6$ matrices at $\alpha = 0.5$. No GPU is required---the entire experiment
operates on 36-parameter matrices via linear algebra.

\paragraph{Interpolation trajectories.}
The three task pairs exhibit qualitatively different merge behavior
(Table~\ref{tab:merge-linearity}). Alpaca--code interpolation is
approximately linear ($R^2 = 0.948$): Fiedler decreases smoothly from
the alpaca endpoint (0.040) to the code endpoint (0.018) with
max residual 0.0035. In contrast, alpaca--math ($R^2 = 0.735$) and
code--math ($R^2 = 0.782$) exhibit non-linear trajectories. The
code--math pair shows a \emph{Fiedler dip} at $\alpha \approx 0.2$:
connectivity drops from 0.018 to 0.015 before recovering to 0.028,
indicating destructive interference at intermediate mixing ratios between
structurally dissimilar adapters.

\begin{table}[t]
\centering
\caption{Bridge interpolation linearity and structure preservation.
Non-linearity correlates with structural dissimilarity between tasks.
Eigenspectrum cosine $> 0.999$ at all $\alpha$ for task merges vs.\
$\sim$0.5 for random baselines.}
\label{tab:merge-linearity}
\small
\begin{tabular}{@{}lccc@{}}
\toprule
Pair & $R^2$ (Fiedler) & Max Residual & Eig.\ $\cos$ @ $\alpha\!=\!0.5$ \\
\midrule
alpaca $\leftrightarrow$ code & 0.948 & 0.0035 & 0.9999 \\
alpaca $\leftrightarrow$ math & 0.735 & 0.0044 & 0.9999 \\
code $\leftrightarrow$ math   & 0.782 & 0.0035 & 1.0000 \\
\midrule
\textit{Random baseline}      & --- & --- & $\sim$0.50 \\
\bottomrule
\end{tabular}
\end{table}

\paragraph{Structure preservation.}
At 10\% contamination ($\alpha = 0.1$), the source adapter's eigenspectrum
is perfectly preserved: cosine similarity $= 1.0000$ for all three pairs.
This holds symmetrically ($\alpha = 0.9$). Small amounts of cross-task
information can be injected into the bridge without disturbing the dominant
task's spectral structure. Even at $\alpha = 0.5$ (equal mixing),
eigenspectrum cosine remains $> 0.999$ for all task pairs. In contrast,
the random baseline produces cosine $\approx 0.5$ at $\alpha = 0.5$:
random noise destroys spectral structure, while cross-task interpolation
preserves it. Bridge merging is structured, not noise.

\paragraph{Per-module merge sensitivity.}
Merge sensitivity (normalized Fiedler change from $\alpha = 0.0$ to $0.5$)
varies by module type. Q-projections are the most merge-sensitive in 2 of
3 pairs (alpaca--code: 0.429; code--math: 0.160), consistent with their
status as the most task-discriminative module (\S\ref{sec:exp-fingerprint}).
The cross-phase correlation between Phase~1A task distance and Phase~2A
merge sensitivity is $r = 0.921$ across the four module types: the modules
that carry the most task identity are the most vulnerable to cross-task
interference. K-proj and V-proj bridges are merge-robust, suggesting a
practical composition strategy: merge K/V bridges freely, interpolate O
bridges cautiously, and maintain task-specific Q bridges.

% ---------------------------------------------------------------------------
\subsection{Overfitting Diagnostic}
\label{sec:exp-overfit}

\textit{This experiment is currently in progress.} We train RhombiLoRA on a
deliberately small split (500 train / 500 validation from Alpaca-cleaned)
for 10,000 steps with checkpoints every 100 steps, tracking Fiedler value
and deviation alongside the train--validation loss gap. The hypothesis is
that bridge spectral properties correlate with overfitting onset ($r > 0.7$)
or exhibit a phase transition at the point where generalization degrades.
Results will be reported in the final version. Training is approximately
21\% complete at the time of writing (estimated 7 hours remaining on an
RTX 4090).
